\chapter{Diferenciación}

\section{La derivada}

\begin{definition}
  $f'(a)=\displaystyle\lim_{h\to0}\frac{f(a+h)-f(a)}{h}$ si el límite existe.
\end{definition}

\begin{proposition}
  Si $f$ es derivable en $a$, entonces es continua en $a$ (lema~\ref{lem:comp}).
\end{proposition}

\begin{remark}
  \begin{align}
    (fg)' &= f'g+fg', \label{eq:leibniz}\\
    (f\circ g)' &= (f'\circ g)\,g'.
  \end{align}
  La ecuación~\eqref{eq:leibniz} es la regla de Leibniz.
\end{remark}

\section{Teorema del valor medio}

\begin{theorem}[del valor medio]
  Si $f$ es continua en $[a,b]$ y derivable en $(a,b)$, existe $c$ con
  $f'(c)=\dfrac{f(b)-f(a)}{b-a}$.
\end{theorem}

\begin{corollary}
  Si $f'=0$ en un intervalo, $f$ es constante.
\end{corollary}
\begin{proof}
  Sean $x<y$ en el intervalo. Por el teorema del valor medio existe $c\in(x,y)$ con
  \[
    f(y)-f(x)=f'(c)\,(y-x)=0.
  \]

  Luego $f(x)=f(y)$ para todo par de puntos, es decir, $f$ es constante.
\end{proof}

\begin{theorem}[Teorema de la función inversa]
  Sea \( W\subset\R^n \) un abierto y \( \phi: W\to \R^n \) una función de clase \( C^1 \). Si para un punto \( P\in W  \) la matriz jacobiana \( J_{\phi}(P) \) es invertible, entonces existen entornos abiertos \( U \) de \( P \) y \( V \) de \( \phi(P) \) tales que la restricción \( \phi: U\to V \) es un difeomorfismo de clase \( C^1 \). Más aún, para todo \( y\in V \) con \( x = \phi^{-1}(y) \), se cumple
  \[
    J_{\phi^{-1}}(y) = [J_{\phi}(x)]^{-1}.
  \]
\end{theorem}